\section{Continuation of solutions in two variables}
\label{cross:section}\label{sec:product}

The local data give a matrix Green kernel when at least one variable lies
in the region where the two structures already agree.  We must continue
this kernel to separated pairs outside that region.  This section proves
the required continuation for elliptic systems acting on separate tensor
indices.  The geometric construction and its scalar analytic foundation
are those of \cite[Section~3]{ScalarCompanion2026}; we reproduce the
argument, including the interpolation and overlap constructions, and
make the extension to systems explicit.  No scalar inverse-problem
uniqueness theorem is used.

\subsection{Separate equations and extension from a cross}

Let $V_i$ be finite-dimensional Hermitian vector spaces.  In a coordinate
domain in $\mathbb R^n$, consider an operator
\[
 \mathsf L_i=-\Delta_{g_i}I_{V_i}
                   +C_i^a\partial_a+D_i,
\]
where $g_i$ is a smooth Riemannian metric and $C_i^a,D_i$ are smooth
endomorphism-valued coefficients.  A \emph{product solution} is a
$V_1\otimes V_2$-valued function $F(x,y)$ satisfying
\[
 (\mathsf L_1^x\otimes I)F=0,
 \qquad (I\otimes\mathsf L_2^y)F=0.
\]
We suppress the identity factors below.  The coefficients of one equation
act only on its indicated tensor index and depend only on its indicated
variable.  In particular, the two equations commute.  Their sum has
scalar real elliptic principal part on the product; ordinary unique
continuation therefore identifies product solutions that agree on a
nonempty open subset of a connected common domain.  Multiplication of
either equation by a positive scalar function does not change its
solutions or any continuation constructed from them.

For matrix Green kernels the first tensor factor is the output fiber.
The second is the conjugate of the input fiber: the equation in $y$ is
obtained by conjugating the coefficients of the second connection
Laplacian.  Equivalently, the adjoint columns of the kernel solve the
second equation.  This convention places the two matrix actions on
separate indices, as required here.

The first construction extends compatible data on
$(D_1\times S_2)\cup(S_1\times D_2)$.  Its mechanism is a basis that is
orthogonal both on a large domain and after restriction to the observed
set.  Common-basis methods have a classical antecedent in separately
harmonic extension \cite{Zeriahi1982}; singular systems also enter
quantitative Runge approximation \cite{RulandSalo2019}.  The interpolation
threshold and the construction below follow the scalar companion.

\begin{lemma}[Extension from a cross]\label{cross:series}
Let $D_i$ be bounded connected coordinate domains, and let
$S_i\Subset D_i$ be nonempty open sets.  Suppose a product solution $F$
is given consistently on
\[
 (D_1\times S_2)\cup(S_1\times D_2),
\]
with the sum of its two $L^2$ norms at most $M$.  Suppose
$O_i\Subset D_i$ are open sets and every $V_i$-valued solution $w$ of
$\mathsf L_iw=0$ on $D_i$ satisfies
\begin{equation}\label{cross:interpolation-input}
 \|w\|_{L^2(O_i)}
 \le C\|w\|_{L^2(S_i)}^{\gamma_i}
          \|w\|_{L^2(D_i)}^{1-\gamma_i},
 \qquad 0<\gamma_i<1,\qquad \gamma_1+\gamma_2>1.
\end{equation}
Then $F$ extends as a product solution to $O_1\times O_2$, agreeing with
the two given pieces on their intersections.  Its $C^m$ norm on every
compact subset of the output product is at most $C_mM$.
These constants are uniform when geometric margins, ellipticity, the
required coefficient bounds, interpolation constants, and a positive
lower bound for $\gamma_1+\gamma_2-1$ are uniform.  A specified output
order requires only finitely many coefficient bounds.
\end{lemma}

\begin{proof}
Let $\mathcal H_1(D_1)$ be the closed subspace of
$L^2(D_1;V_1)$ consisting of distributional $\mathsf L_1$-solutions.
Restriction $R:\mathcal H_1(D_1)\to L^2(S_1;V_1)$ is compact by interior
elliptic estimates and is injective by unique continuation.  The spectral
theorem for $R^*R$ gives an orthonormal basis $(w_j)$ of
$\mathcal H_1(D_1)$ such that
\begin{equation}\label{cross:singular-values}
 (w_i,w_j)_{L^2(S_1)}=\mu_j^2\delta_{ij},\qquad 0<\mu_j\le1.
\end{equation}
For every $N$ there is a uniform constant $C_N$ with
\begin{equation}\label{cross:singular-decay}
 \mu_j\le C_N(1+j)^{-N}.
\end{equation}
Indeed, multiply a solution by a cutoff equal to one near
$\overline{S_1}$ and compactly supported in $D_1$.  Interior estimates
bound its $H^k$ norm by its $L^2(D_1)$ norm.  Fourier truncation in a
containing cube, component by component, gives an approximation of rank
$O(A^n)$ and error $O(A^{-k})$.  The minimax characterization of singular
values yields $\mu_j\le C_kj^{-k/n}$.  The fixed fiber dimension changes
only the constants.

For $y\in S_2$, expansion in the first variable gives
\[
 F(x,y)=\sum_j w_j(x)\otimes b_j(y),\qquad
 \|b_j\|_{L^2(S_2;V_2)}\le M.
\]
Contracting the first tensor index against $\overline{w_j}$ extends the
coefficients to $D_2$:
\begin{equation}\label{cross:coefficient-extension}
 b_j(y)=\mu_j^{-2}\int_{S_1}
                  \sum_\alpha\overline{w_{j,\alpha}(x)}F_{\alpha}(x,y)\,dx,
 \qquad \|b_j\|_{L^2(D_2;V_2)}\le M\mu_j^{-1}.
\end{equation}
Here $F_\alpha(x,y)\in V_2$ are the first-index components in an
orthonormal basis of $V_1$; the displayed contraction leaves a vector
in $V_2$.  Because the
second system acts on that remaining index, $\mathsf L_2b_j=0$.
Restriction orthogonality proves that this formula agrees with the
original coefficient on $S_2$.  The assumed interpolation gives
\[
 \|w_j\|_{L^2(O_1)}\le C\mu_j^{\gamma_1},\qquad
 \|b_j\|_{L^2(O_2)}\le CM\mu_j^{\gamma_2-1}.
\]
Thus the tensor series converges absolutely in $L^2(O_1\times O_2)$:
the norm of its $j$th term is at most
$CM\mu_j^{\gamma_1+\gamma_2-1}$, and
\eqref{cross:singular-decay} makes these bounds summable.  The limit
solves both equations distributionally.  Interior estimates for their
elliptic sum give smoothness and the claimed bounds.

Agreement on the second arm requires a completeness check.  Take
$z\in L^2(S_1;V_1)$ orthogonal to the closure of
$R\mathcal H_1(D_1)$.  The $V_2$-valued function
\[
 y\longmapsto\int_{S_1}\sum_\alpha\overline{z_\alpha(x)}F_\alpha(x,y)\,dx
\]
solves the second system and vanishes on $S_2$, hence vanishes on
connected $D_2$.  Every second-arm slice therefore belongs to the
closed restricted range.  Its expansion in the orthonormal basis
$w_j/\mu_j$ is complete, with coefficients $\mu_jb_j(y)$.
Consequently the series equals $F$ on $S_1\times O_2$; it converges
there by the positive exponent $\gamma_2$.  On $O_1\times S_2$ it is the
original expansion, converging with exponent $\gamma_1$.  These establish
both required agreements.  Only upper bounds for $\mu_j$ have been
used, so uniformity does not require continuous choices of singular
bases or lower bounds on singular values.
\end{proof}

The remaining issue is geometric: we must place two observed sets so
that the interpolation exponents add to more than one.  Strict Carleman
weights will supply those exponents from a curvature condition on the
boundary of the region where the product solution is known.

\subsection{A geometric criterion for local extension}

We next build the interpolation estimates needed in
Lemma~\ref{cross:series}.  For a metric $g$, call a smooth real function
$\theta$ a \emph{strict weight} at a point where $d\theta\ne0$ if
\begin{equation}\label{cross:strict-weight}
 \operatorname{Hess}_g\theta(e,e)
  +\operatorname{Hess}_g\theta(v,v)>0
 \quad\text{for all }v\perp e,\ |v|_g=1,
 \qquad e=\frac{\nabla^g\theta}{|\nabla^g\theta|_g}.
\end{equation}
For the system $\mathsf L=-\Delta_gI+C^a\partial_a+D$, the strict
Carleman estimate gives, after shrinking the coordinate neighborhood,
\begin{equation}\label{cross:carleman}
 \tau^3\|e^{\tau\theta}w\|_{L^2}^2
 +\tau\|\nabla(e^{\tau\theta}w)\|_{L^2}^2
 \le C\|e^{\tau\theta}\mathsf Lw\|_{L^2}^2,
 \qquad \tau\ge\tau_0,
\end{equation}
for compactly supported vector-valued $w$.  All norms sum over the
fiber components.  The scalar form is the elliptic strict-weight estimate
of the classical
Carleman theory; see \cite[Chapter~XXVIII]{HormanderIV2009} and
\cite[Definition~8.3 and Theorem~9(a)]{KochTataru2005}.
We recall its energy proof to record the system dependence.  First take
one scalar component and $\mathsf L=-\Delta_g$,
write $v=e^{\tau\theta}w$ and
\[
 e^{\tau\theta}\mathsf Le^{-\tau\theta}=A_\tau+B_\tau,\qquad
 A_\tau=-\Delta_g-\tau^2|d\theta|_g^2,\quad
 B_\tau=\tau(2\nabla^g\theta\cdot\nabla+\Delta_g\theta).
\]
Here $A_\tau$ is self-adjoint and $B_\tau$ is skew-adjoint for
metric volume.  Integration by parts gives
\begin{align*}
 \|(A_\tau+B_\tau)v\|^2
 &=\|A_\tau v\|^2+\|B_\tau v\|^2
       +([A_\tau,B_\tau]v,v),\\
 ([A_\tau,B_\tau]v,v)
 &=4\tau\int\operatorname{Hess}_g\theta(\nabla v,\nabla\overline v)
   +4\tau^3\int\operatorname{Hess}_g\theta
                 (\nabla\theta,\nabla\theta)|v|^2
   -\tau\int(\Delta_g^2\theta)|v|^2.
\end{align*}
Choose a real constant $m$ strictly between the negative of the least
unit tangential Hessian value and the unit normal Hessian value.
After shrinking the patch, these inequalities have a positive margin.
Add $4m\tau(A_\tau v,v)$ to the commutator form.  Its gradient
coefficient becomes $4\tau(\operatorname{Hess}_g\theta+mg)$,
and its leading zero-order coefficient is
\[
 4\tau^3|d\theta|_g^2
       \big(\operatorname{Hess}_g\theta(e,e)-m\big)>c\tau^3.
\]
The tangential gradient form is positive.  Its mixed and possibly
negative normal terms cost at most
$C\tau\|\nabla_e v\|^2$, and
\[
 \|\nabla_e v\|^2
 \le C\tau^{-2}\|B_\tau v\|^2+C\|v\|^2.
\]
Finally
$|4m\tau(A_\tau v,v)|\le\tfrac12\|A_\tau v\|^2+
C\tau^2\|v\|^2$.
For large $\tau$, the squared parts absorb the normal derivative
cost and the positive zero-order term absorbs the remaining errors.
This yields the scalar estimate.  Summing it over components gives the
estimate for $-\Delta_gI$.  The conjugated matrix lower-order terms have
norm at most $C(\|\nabla v\|+\tau\|v\|)$, so their squared norm is
absorbed by the $\tau$ and $\tau^3$ terms after increasing $\tau_0$.
This proves \eqref{cross:carleman} for the stated systems.  The constants
are uniform on compact sets of strict weights, with fixed positive
strictness margins and bounded matrix coefficients.  The scalar cutoff
commutators used below act componentwise, and interior estimates for
these systems have the same orders as in the scalar case.

The existence criterion below is the geometric criterion of
\cite[Section~3, ``A geometric criterion for local extension'']{ScalarCompanion2026}.
Its geometric proof depends only on the principal metrics.  The system
estimate just proved and Lemma~\ref{cross:series} provide the analytic
steps in the present setting.  All gradients, lengths, orthogonality
relations, and Hessians in its statement use the principal metrics
$g_1,g_2$ and their product connection.

\begin{proposition}[Product extension criterion]\label{cross:criterion}
Let $\rho$ be smooth near $z_0=(x_0,y_0)$, with $\rho(z_0)=0$ and
$d_x\rho,d_y\rho\ne0$.  Put
\[
 a_i=\frac{\nabla_i\rho}{|\nabla_i\rho|^2},\qquad H=\operatorname{Hess}\rho.
\]
Suppose that at $z_0$
\begin{equation}\label{cross:hessian-test}
 H((a_1,-a_2),(a_1,-a_2))
 +H((v_1,v_2),(v_1,v_2))>0
 \quad
 \left(v_i\perp a_i,\ |v_i|=|a_i|\right).
\end{equation}
Every product solution given on $\{\rho>0\}$ near $z_0$ extends to
a neighborhood of $z_0$, agreeing on a smaller part of that side.
In fact, the construction uses only two compact product sets in
$\{\rho>\varepsilon\}$ for some $\varepsilon>0$.
The neighborhoods and estimates on smaller neighborhoods are uniform
under small smooth perturbations preserving the strict hypotheses,
provided the data on those compact sets are bounded.
\end{proposition}

\begin{proof}
We first split the product Hessian condition into a strict weight in
each factor. Their sum will lie below a defining function for the given
side, with a quadratic gap. That gap places two intersecting product
sets in the given side; the weights then yield interpolation exponents
greater than one half, allowing Lemma~\ref{cross:series} to apply.

\paragraph{Splitting the weight between the factors.}
We seek symmetric forms $P_i$ satisfying
\[
 P_i(a_i,a_i)+P_i(v_i,v_i)>0,
 \qquad H+K\,d\rho^{\otimes2}>P_1\oplus P_2
\]
for some $K>0$, with the vectors $v_i$ as in
\eqref{cross:hessian-test}. After division by $|a_i|^2$, the first
inequality is the strict-weight condition for a function with gradient
$\nabla_i\rho$ and Hessian $P_i$. The second inequality will give the
quadratic gap. To obtain these forms by convex separation, the dual
tests are positive semidefinite forms $D$ whose diagonal blocks are
\[
 d_i a_i^{\otimes2}+S_i,\qquad
 d_i\ge0,\quad S_i\ge0\text{ on }a_i^\perp,
 \quad \operatorname{tr} S_i=d_i|a_i|^2.
\]
This follows by separating the open convex cone consisting of positive
definite forms plus the allowed $P_1\oplus P_2$; the individual dual
cones are generated by $a_i^{\otimes2}+v_i^{\otimes2}$.
Normalize $\operatorname{tr} D=1$.  To obtain a suitable $K$, it suffices by
compactness to prove $H:D>0$ on tests with $D\,d\rho=0$.

Represent such a $D$ as a covariance matrix.  The normal components
measured by the two partial covectors are opposite.  The diagonal-block
condition makes each normal component uncorrelated with its own
tangential component, hence with both tangential components.
Their variances coincide, giving $d_1=d_2=d>0$ and
\[
 D=d(a_1,-a_2)^{\otimes2}+D_\perp,
 \qquad \operatorname{tr}(D_\perp)_{ii}=d|a_i|^2.
\]
The case $d=0$ would force $D=0$.  Among nonnegative tangential forms
with these two fixed traces, every extreme point has rank one:
on an image of rank $r\ge2$, a nonzero symmetric perturbation preserves
the two traces because $r(r+1)/2>2$, and small perturbations of both
signs preserve positivity.  The rank-one tests are precisely those
in \eqref{cross:hessian-test}, multiplied by $d$.
The asserted positivity follows.  Compactness of the normalized tests
then supplies $K$, as required.

Prescribe partial covectors $d_x\rho,d_y\rho$ and Hessians $P_i$
for smooth functions
$\theta_i$ vanishing at the respective points.  In small coordinates
centered there, Taylor expansion gives
\begin{equation}\label{cross:quadratic-gap}
 \theta_1(x)+\theta_2(y)
 \le \widetilde\rho(x,y)-c(|x|^2+|y|^2),
 \qquad \widetilde\rho=\rho+K\rho^2/2,
\end{equation}
with $c>0$.  Each $\theta_i$ satisfies
\eqref{cross:strict-weight}.  The positive side of $\widetilde\rho$
is the positive side of $\rho$ in this neighborhood.

\paragraph{Placing the cross and obtaining interpolation.}
Use coordinates $(s_i,z_i)$ with $s_i=\theta_i$. Fix a small
$\kappa>0$ so that $\phi(s_i)=s_i-\kappa s_i^2$ remains a strict weight.
Choose successively a small lateral radius $R$, a height
$l\ll cR^2$, and $e\ll\kappa l^2$.  Take
\[
 D_i=\{-l<s_i<l+4e,\ |z_i|<R\}.
\]
Use a cutoff equal to one on the inner lateral half and for
$-l+3e<s_i<l+2e$, supported in the cylinder and in
$-l+2e<s_i<l+3e$.
Let $S_i\Subset D_i$ include neighborhoods of its top and lateral
derivative supports, where respectively
\[
 s_i>l+e/2\quad\text{or}\quad |z_i|>R/3,
\]
and an inner top slice near $s_i=l+5e/4$.
The two closed product sets $\overline{D_1\times S_2}$ and
$\overline{S_1\times D_2}$ lie strictly in the given side.
For a top slice, the sum of the two $s$ coordinates is positive.
For a lateral slice, it is greater than $-2l$, which is dominated
by the quadratic gap in \eqref{cross:quadratic-gap}.
Making the sets slightly smaller if necessary gives a common
$\varepsilon>0$ with $\rho>\varepsilon$ on their closures.

For an $\mathsf L_i$-solution normalized by $\|w\|_{L^2(D_i)}=1$, apply
\eqref{cross:carleman} to this cutoff times $w$.  Interior estimates
bound the bottom derivative terms using the large norm and the weight
$\phi(-l+3e)$.  The remaining derivative terms use
$\|w\|_{L^2(S_i)}$ and a weight at most $\phi(l+3e)$.
On a smaller inner cylinder beginning at $s_i=-e$, the weight is at
least $\phi(-e)$.  Balancing the two exponentials yields
\eqref{cross:interpolation-input} with
\begin{equation}\label{cross:exponent}
 \gamma_i=
 \frac{\phi(-e)-\phi(-l+3e)}{
       \phi(l+3e)-\phi(-l+3e)}>\frac12.
\end{equation}
Indeed at $e=0$ the quotient is $1/2+\kappa l/2$.
Values of the balancing parameter below $\tau_0$ are covered by
increasing the constant.  The output cylinder may be chosen to include
a connected tube from the origin to the indicated top slice.

\paragraph{Constructing and identifying the extension.}
Lemma~\ref{cross:series} now gives a product solution on a neighborhood
of $z_0$. It agrees with the old solution there on the positive side:
from a sufficiently near positive-side point, increase $s_2$ into
the top slice.  Since $\partial_{s_2}\rho>0$ locally, the segment stays
on that side and in the larger constructed cylinder.  Agreement on
the slice and unique continuation along a neighborhood of the segment
prove agreement at the original point.
All choices have strict margins.  Keeping their finitely many compact
sets and shrinking their output neighborhoods once proves the stated
perturbation uniformity by \eqref{cross:carleman},
Lemma~\ref{cross:series} and interior elliptic estimates.
\end{proof}

\subsection{Gluing along a compact family of surfaces}

Proposition~\ref{cross:criterion} provides a germ across one surface
point.  The following formulation records the geometry needed to
combine these germs.  In particular, agreement is proved on specified
overlap components, rather than inferred from the existence of a cover.

\begin{lemma}[Propagation through a compact cylinder]
\label{cross:propagation}
Suppose a region in a product patch has smooth coordinates $(z,a)$
near $B\times[a_-,a_+]$, where $B$ is a compact base, possibly with
boundary or corners.  A single product solution is given on an open set $\mathcal S$ containing
neighborhoods of the top and lateral boundary.
Assume that $\mathcal S$ is upward closed in $a$ at fixed $z$.  If every level $a=\text{constant}$
satisfies Proposition~\ref{cross:criterion} outside $\mathcal S$, with
the positive side pointing toward increasing $a$, the solution extends
to a neighborhood of the cylinder, agreeing on $\mathcal S$.
For fixed coordinates, compact sets and strict reference inequalities,
the extension has uniform compact-subset bounds under small smooth
perturbations, in terms of bounds on finitely many compact subsets of
the initially given region.
\end{lemma}

\begin{proof}
Starting from the top, let $a_*$ be the infimum of levels down to which
an extension agreeing with the seed has been obtained.  In the uniform
version include uniform bounds in this property.  At a non-seed point
of $B\times\{a_*\}$, Proposition~\ref{cross:criterion} uses compact
data strictly above that level.  Take a finite cover of the level by
such output neighborhoods and seed neighborhoods.  Their data are
contained above $a_*+\eta$ for one $\eta>0$, so an already reached
level supplies them with the needed bounds.

Shrink the outputs in $(z,a)$ coordinates to boxes
$B_i\times(a_*-d_i,a_*+d_i)$.
Every connected component of an overlap contains vertical segments
to a common height greater than $a_*$.  Both germs equal the preceding
extension there.  Unique continuation for $\mathsf L_1^x+\mathsf L_2^y$ makes them
identical on that overlap component.  The same argument gives agreement
with $\mathcal S$, because a vertical segment moving upward from a
seed point stays in the seed.  Boundary neighborhoods use the given
solution.  A finite subcover has a positive minimum output width and
therefore passes the putative last level, or includes the endpoint
$a_-$.  This proves continuation on the whole cylinder.
The same finite constructions retain all strict margins under small
perturbations.  Their Carleman, series and interior estimates prove
the uniform assertion.  Data bounded up to a limiting surface are
never required: each local construction uses compact sets strictly
above it.
\end{proof}

\subsection{Initialization at a fixed distance from the diagonal}

We now return to the connection Laplacians.  In common coordinate
patches centered at $0$, let $(g_j,A_j)$ be smooth metrics and unitary
connection matrices obtained from interior patches of compact Dirichlet
manifolds.  Use fixed local unitary frames, and let $G_j$ denote their
matrix Dirichlet Green kernels, normalized with metric volume.
Suppose that both coefficient pairs agree on an open set $W$ containing
the lower side of a smooth hypersurface through $0$, and that
$G_1(x,y)=G_2(x,y)$ for distinct $x,y\in W$.  After a common linear
coordinate change, assume
\[
 g_1(0)=g_2(0)=I_n,\qquad
 \{x_n<q(x')\}\subset W\text{ near }0,\qquad
 q(0)=0,\quad d q(0)=0,
\]
for a smooth function $q$.  No regularity of the rest of $\partial W$
is assumed.
For a small spatial dilation parameter $\lambda>0$, put
\begin{equation}\label{cross:dilation}
 \begin{gathered}
 g_{j,\lambda}(z)=g_j(\lambda z),\qquad
 A_{j,\lambda}(z)=\lambda A_j(\lambda z),\qquad
 W_\lambda=\lambda^{-1}W,\\
 G_{j,\lambda}(x,y)=\lambda^{n-2}G_j(\lambda x,\lambda y).
 \end{gathered}
\end{equation}
On each fixed bounded patch the metrics tend smoothly to the Euclidean
metric.  The corresponding dilated systems have uniformly bounded
smooth lower-order coefficients on every fixed bounded set, and
$W_\lambda$ contains every fixed compact subset of $\{z_n<0\}$ for
sufficiently small $\lambda$.
Initially define, off the diagonal,
\begin{equation}\label{cross:mixed-data}
 \mathcal G_\lambda(x,y)=
 \begin{cases}
 G_{1,\lambda}(x,y),&y\in W_\lambda,\\
 G_{2,\lambda}(x,y),&x\in W_\lambda.
 \end{cases}
\end{equation}
The pieces agree on their overlap.  They solve the first connection
Laplacian in $x$ and the conjugate second connection Laplacian on the
second index in $y$, in the convention at the beginning of the section.
Indeed, on the observed factor its two coefficient systems agree.
Thus \eqref{cross:mixed-data} is a product solution of the systems
already treated.

\begin{proposition}[Fixed-separation initialization]\label{cross:initial}
In this setting, let $n\ge3$ and fix $B>0$, $d>0$ and $\eta>0$.
For all sufficiently small $\lambda$, the data
\eqref{cross:mixed-data} have a product-solution continuation to a
neighborhood of
\[
 \{(x,y):|x|,|y|\le B,\ |x-y|\ge d\}.
\]
Every fixed $C^m$ norm on a smaller neighborhood is bounded uniformly
in $\lambda$.  It agrees with $G_{1,\lambda}$ when $y\in W_\lambda$
and with $G_{2,\lambda}$ when $x\in W_\lambda$, on the corresponding
overlaps.  In particular these agreements hold on the fixed truncated
lower strips $y_n<-\eta$ and $x_n<-\eta$, respectively.  The smallness threshold for $\lambda$ and
the constants may depend on $B,d,\eta,m$.
\end{proposition}

\begin{proof}
For the final comparison with the full original cross, fix at the outset
\[
 \widehat B=B+\eta+4,\qquad
 \widehat d=\min(d/2,1/4).
\]
We construct the continuation for these enlarged location bounds and
smaller separations, and then restrict to the stated target.
We first cross the limiting plane, then deform variable-radius tubes
to reach those separations.  This follows the two-stage
construction of \cite[Section~3, ``Fixed-separation initialization'']{ScalarCompanion2026}.
We first choose the Euclidean geometric ratios, including the opening of
the region reached in the first stage.  We next choose the entire tube
cylinder and a bounded positive-separation range containing it.  The
first-stage constructions are then made on that range.  Only after all
these fixed sets have been chosen do we decrease $\lambda$.

\paragraph{Crossing the plane and obtaining a common seed.}
We first extend slightly across the plane in one variable while keeping
the other at a comparable positive separation. The first interpolation
exponent can be made close to one; this compensates for the fixed
positive exponent obtained by propagation in the separated second
variable.

For the Euclidean equations, the following choices are invariant under
translations parallel to the plane and under a common change of spatial
scale.  We make the geometric choices on a unit-scale template.  They
therefore give an opening constant $c_0>0$ independent of the bounded
range on which the construction is subsequently used.  On any fixed
compact range of centers and scales $0<s_{\min}\le s\le s_{\max}$,
the same strict weights and geometric fractions work for the dilated
systems once $\lambda$ is sufficiently small.  The lower-order
coefficients are bounded uniformly there (and converge to zero under
the connection dilation).

Fix a point $x_0$ of height zero and a separation scale $s>0$.
In the second factor take
\[
 D_2=\{s/2<|y-x_0|<2s\}
\]
and an observed ball $S_2$ about $x_0-se_n$.
In the first take a ball of radius $3\vartheta_0s$ centered at
$x_0-\vartheta_0se_n$, with $\vartheta_0>0$ fixed small enough that the two
large domains are separated.  Its observed ball has radius
$(1-\alpha)\vartheta_0s$, with $\alpha>0$ to be chosen.
These observed balls lie strictly in the lower half-space.  On a fixed
compact scale range their heights are bounded above by a common negative
number, since $s\ge s_{\min}>0$.  Thus, after a finite covering, one
choice of small $\lambda$ places every observed ball in $W_\lambda$
with room.  No neighborhood of the whole limiting plane is assumed to
belong to $W_\lambda$.

There is a two-constants estimate from $S_2$ to a slightly enlarged
closed annulus $4s/5\le|y-x_0|\le6s/5$, with an exponent
$\gamma_2>0$ independent of $\alpha$.  Here is a direct way to obtain
it with room for perturbations.  On a Euclidean annulus the weight
$|z-z_*|^{-1}$ is strict: its radial Hessian is
$2|z-z_*|^{-3}$ and each tangential Hessian value is
$-|z-z_*|^{-3}$.  Radial cutoffs in
\eqref{cross:carleman}, followed by exponential balancing, propagate
smallness from an inner ball to an intermediate ball using the norm
on an outer ball.  A finite chain of overlapping balls with enlarged
balls inside $D_2$ connects $S_2$ to the target annulus.  The annulus
is connected since $n\ge3$.  Multiplying the finitely many positive
exponents gives the asserted $\gamma_2$.  The same weights and margins
work for sufficiently small perturbations of the Euclidean metric.

In the first factor use this radial argument centered at
$x_0-\vartheta_0se_n$.  Keep the outer cutoff a fixed distance beyond radius
$\vartheta_0s$, and put the inner cutoff just inside radius
$(1-\alpha)\vartheta_0s$.  Take a concentric target ball of radius
$(\vartheta_0+c)s$, with $c>0$ small; it contains the observed ball and
a ball of radius $cs$ about $x_0$.
To see the resulting exponent, let $R_{\rm in}$ be an inner cutoff
radius, chosen just inside the observed radius, let
$R_{\rm tar}=(\vartheta_0+c)s$, and let $R_{\rm out}>R_{\rm tar}$ be
a fixed starting radius for the outer derivative support.  With
$\theta(R)=R^{-1}$, exponential balancing gives
\[
 \gamma_1=
 \frac{\theta(R_{\rm tar})-\theta(R_{\rm out})}
      {\theta(R_{\rm in})-\theta(R_{\rm out})}.
\]
The inner core is bounded directly by the observed norm.  As $\alpha$
and $c$ decrease, choose the inner cutoff so that both
$R_{\rm in}$ and $R_{\rm tar}$ approach $\vartheta_0s$, while
$R_{\rm out}$ stays a fixed distance beyond that radius.  Thus
$\gamma_1\to1$.  Choose these constants so that
$\gamma_1+\gamma_2>1$.
Lemma~\ref{cross:series} extends the kernel to a small ball about
$x_0$ times the target annulus.  On these fixed separated sets the
rescaled Green kernels and their derivatives are uniformly bounded:
the same bounds hold entrywise for the smooth systems used here.  To see
this, an interior inverse parametrix has order $-2$ and frequency pieces
bounded by $C_N2^{j(n-2)}(1+2^j|x-y|)^{-N}$, up to a smooth kernel.
Summing gives the local bound $C|x-y|^{2-n}$ for $n\ge3$.
The normalization in \eqref{cross:dilation} preserves this bound;
interior estimates for the uniformly elliptic dilated equations give all
fixed derivative bounds away from the pole.

The extensions agree with both lower pieces on their actual comparison
domains.  For $y\in S_2$, the series equals $G_{1,\lambda}$ for all
$x$ in its first output ball.  Fix any such $x$ in a truncated lower
half-space contained in $W_\lambda$.  On $S_2$ both variables then lie
in $W_\lambda$, so this value also equals $G_{2,\lambda}$.  Unique
continuation for the second system on the connected output annulus
proves agreement with $G_{2,\lambda}$ throughout that annulus.  In the
other direction, fix $y$ in the output annulus and in $W_\lambda$.
On the first observed ball $S_1$, agreement with the second arm gives
$G_{2,\lambda}=G_{1,\lambda}$.  Unique continuation for the first
system on its connected output ball gives agreement with
$G_{1,\lambda}$ there.  Each comparison stays a positive distance from
the diagonal.

We make the overlap comparison precise because the two variables may
play different roles in two local constructions.  Inside the output just
obtained choose a larger comparison product of the form
\[
 P(x_0,s)=B(x_0,\kappa s)\times
 \{(1-\zeta)s<|y-x_0|<(1+\zeta)s\},
\]
where $\kappa,\zeta>0$ are fixed geometric fractions.  The first ball
is contained in the actual concentric first output ball, and the annulus
is contained in the actual second output.  Keep the full original
outputs for the preceding lower-piece comparisons.  For coverage retain
only
\[
 P'(x_0,s)=B(x_0,\epsilon_1s)\times
 \{(1-\zeta')s<|y-x_0|<(1+\zeta')s\},
 \qquad 0<\epsilon_1\ll\kappa,\quad 0<\zeta'<\zeta.
\]
The transposed products are used when $y$ is the near-plane variable.
If a point belongs to two retained products, their scales are comparable
to each other and to $|x-y|$, with fixed constants.  Choose $\epsilon_1$
small enough, relative to $\kappa$ and $\zeta-\zeta'$, that a downward
translation of size at most a fixed multiple of $\epsilon_1$ times
either scale stays inside each larger comparison product whenever it
acts on a retained near-plane variable.  This follows directly from
the triangle inequality for the ball and annular distances.  On the
chosen compact scale range take $\eta_0>0$ much smaller than the
remaining margins times $s_{\min}$, and with $\eta_0\le\eta$.

For products of the same orientation, lower only their common near-plane
variable by
\[
 h=\max(x_n,0)+2\eta_0
\]
(or the analogous expression in $y$).  The other variable remains
fixed.  For products of opposite orientations, both $x$ and $y$ lie
in retained near-plane balls; lower both by
\[
 h=\max(x_n,y_n,0)+2\eta_0.
\]
The scale comparability and the preceding choices keep the entire path
in the intersection of the two larger comparison products.  Its
endpoint has the near-plane variable below $-\eta_0$ in the first case,
and both variables below $-\eta_0$ in the second.  For small enough
$\lambda$ these endpoint neighborhoods lie in the original lower cross,
where the preceding comparisons identify both germs.  Unique
continuation for the elliptic sum on the component of the larger
intersection containing the path identifies the germs at the starting
point.  The path need not remain in the retained products.  This proves
agreement on every retained overlap without a connectedness assumption
on that overlap.

Apply the same construction with the variable roles and their tensor
factors interchanged, and then return the values to the original tensor
order.  The resulting germs solve the same two equations as the first
orientation.  The overlap comparison therefore applies to them and
produces, for some fixed $c_0>0$, a single-valued continuation in
\begin{equation}\label{cross:coarse-region}
 \min(x_n,y_n)<c_0|x-y|,
\end{equation}
on every fixed compact range of bounded locations and positive
separations.  The estimates are uniform in $\lambda$ on that range.
For coverage near the plane, take $x_0=(x',0)$ and
$s=|y-x_0|$ when the first variable has the smaller nonnegative height;
then $x$ is in the retained first ball and $y$ in the retained annulus
if $c_0$ is small; for instance
$c_0/(1-c_0)<\epsilon_1$ guarantees the required ball inclusion.
Points already below the plane are covered by the
same construction near it or by the given cross farther below it.
The preceding overlap argument makes these choices compatible.

\paragraph{Sweeping tubes inward from the common seed.}
The region \eqref{cross:coarse-region} is now the seed for the second
stage. At fixed center $y$ and direction $\omega$, we move $x$ along the
ray from $y$: large radii lie in the seed, and decreasing the radius
will reach the target pairs. Choose $M>\widehat B+3$ and introduce
\begin{equation}\label{cross:tubes}
 t=M+y_n+e_0|y'|^2,\qquad r_a(y)=a t^{1+\sigma},\qquad
 x=y+r_a(y)\omega,\quad |\omega|=1,
\end{equation}
where $\sigma,e_0>0$ will be small.  The base is
\[
 y_n\ge-\widehat B-1,\qquad t\le T_0,\qquad \omega\in\mathbb S^{n-1},
\]
and $a$ decreases through a fixed interval $[a_{\min},a_{\max}]$.
The term $e_0|y'|^2$ makes this base compact. The parameter $a$ decreases
while $(y,\omega)$ remains fixed, as required by the compact-cylinder
propagation lemma.
Put $m=M-\widehat B-1>2$.  Reserve the bounds
$0<\sigma\le1$, $e_0\le(1+\widehat B^2)^{-1}$, and
$e_0T_0\le1/4$.  Once $\sigma$ is chosen below, take
\[
 0<a_{\min}<
 \frac{\widehat d}{2(M+\widehat B+1)^{1+\sigma}},
 \qquad
 a_{\max}>\max\left(a_{\min},\frac{2}{c_0m^\sigma}\right),
\]
and then choose
\[
 T_0>2(M+\widehat B+1),\qquad
 c_0a_{\min}T_0^\sigma>2.
\]
For every target center, $t\le M+\widehat B+1$; hence its target
separations exceed $r_{a_{\min}}(y)$.  On the initial level,
$c_0r_{a_{\max}}>2t>y_n$, so this level lies in
\eqref{cross:coarse-region}.  The bottom base boundary has $y_n<0$,
and on the face $t=T_0$ we have $c_0r_a>2t>y_n$ for all allowed $a$.
Thus the initial level and both base boundary faces are supplied by
the seed.  The power $1+\sigma>1$ is what permits this last inequality
at the top face.
This known region is upward closed in $a$: since
$\min(x_n,y_n)=y_n+r\min(\omega_n,0)$, its defining
inequality is $y_n<(c_0-\min(\omega_n,0))r$, whose radius coefficient
is positive.

It remains to check the strict extension criterion outside that region.
At the Euclidean metrics use
\[
 \rho=\tfrac12(|x-y|^2-r_a(y)^2),\qquad b=\nabla r_a.
\]
Away from $\omega+b=0$, multiply the vectors in
\eqref{cross:hessian-test} by the common positive factor $r_a$ and
write their real and imaginary parts as
\begin{equation}\label{cross:null-vectors}
 \begin{aligned}
 U_1&=\omega+iv,& |v|=1,\quad v\perp\omega,\\
 U_2&=\frac{\omega+b}{|\omega+b|^2}+iw,
 &w\perp(\omega+b),\quad |w|=|\omega+b|^{-1}.
 \end{aligned}
\end{equation}
The identity $\omega\cdot(U_1-U_2)=b\cdot U_2$ cancels the normal
part of the distance Hessian.  Consequently the test is
\begin{equation}\label{cross:tube-test}
 \big|\pi_{\omega^\perp}(U_1-U_2)\big|^2
 -(r_a\partial^2r_a)(U_2,\overline{U_2}).
\end{equation}
For bounded $b$, outside fixed small neighborhoods of
$b=-\omega$ and $b=-2\omega$, the first term is at least $c|b|^2$.
To see this, a zero forces the real part of $U_2$ to be parallel to
$\omega$, and equality of the projected imaginary lengths gives
$|\omega+b|=1$.  Thus $b=0$ or $b=-2\omega$.
Near $b=0$, the real projected difference controls
$|\pi_{\omega^\perp}b|$ to first order; also
\[
 |\pi_{\omega^\perp}w|
 =1-\omega\cdot b+O(|b|^2),
\]
so the imaginary projected difference controls $|\omega\cdot b|$
to first order.  This proves the quadratic lower bound near zero;
compactness proves it on the remaining range.

Outside \eqref{cross:coarse-region}, $y_n\ge c_0r_a$, so
$r_a/t\le1/c_0$ and
\[
 b=(1+\sigma)(r_a/t)\nabla t
\]
has $|b|\le4/c_0$ under the reserved bounds.  This fixes a compact
range for $b$ before $\sigma$, the $a$-interval, $T_0$, and $e_0$
receive their final values.  When $\nabla t$ is close to $e_n$, the two excluded
neighborhoods already lie in \eqref{cross:coarse-region}.  Indeed,
\[
 \frac{x_n}{r_a}
 \le \frac{(1+\sigma)|\nabla t|}{|b|}+\omega_n,
\]
which is arbitrarily small near $b=-\omega$ and negative near
$b=-2\omega$ as $\sigma$ and the gradient error tend to zero.
On the complement $|U_2|$ is bounded and differentiation of
\eqref{cross:tubes} gives
\begin{equation}\label{cross:tube-error}
 r_a|\partial^2r_a|
 \le C(\sigma+T_0e_0)|b|^2.
\end{equation}
Here the choices can be made in a definite order.  First choose a
neighborhood radius $\delta_0>0$ for $b=-\omega$ and $b=-2\omega$,
in terms of $c_0$, and upper bounds on $\sigma$ and
$|\nabla t-e_n|$ so that their closures satisfy $x_n/r_a<c_0/2$.
On the complementary bounded $b$-range let $c>0$ be the constant in the
quadratic lower bound.  There $|U_2|^2\le2/\delta_0^2$.
Direct differentiation sharpens \eqref{cross:tube-error} to
\[
 \frac{r_a|\partial^2r_a|}{|b|^2}
 \le\sigma+2e_0T_0.
\]
Choose $\sigma>0$ small enough that its contribution to this bound is
less than $c\delta_0^2/8$.  Choose the $a$-interval and $T_0$ as above,
and finally take $e_0>0$ sufficiently small that all reserved bounds hold,
$|\nabla t-e_n|\le2\sqrt{e_0T_0}$ has the required smallness, and
$\sigma+2e_0T_0<c\delta_0^2/4$.  The second term of
\eqref{cross:tube-test} is then at most $(c/2)|b|^2$ in absolute value.
The whole test is bounded below by
\[
 \frac c2|b|^2
 \ge\frac c2(1+\sigma)^2a_{\min}^2m^{2\sigma}>0
\]
where continuation is needed.  Also the two partial gradients of
$\rho$ have lengths at least $r_a$ and $\delta_0r_a$, respectively.
This supplies fixed strict margins before any metric perturbation.

All these sets are bounded and stay at separation at least
$a_{\min}(M-\widehat B-1)^{1+\sigma}>0$.  Smooth convergence of the metrics
therefore preserves the strict tests for sufficiently small $\lambda$.
At this point the tube parameters and its compact range have all been
fixed.  Apply the first-stage construction on a slightly larger bounded
range of locations and separations containing this cylinder, and then
choose $\lambda$ small enough for both its data and the strict tube
tests.  Lemma~\ref{cross:propagation} continues the kernel through the
cylinder, with uniform estimates and agreement on the seed.  A target pair has
$a=|x-y|/t^{1+\sigma}\ge a_{\min}$; if $a>a_{\max}$ it is already in
the coarse region.  Hence every target pair is covered.

It remains to identify this extension with both pieces of the full
cross \eqref{cross:mixed-data}, including parts of $W_\lambda$ not in
a truncated half-space.  For $|y|\le B$ put
\[
 D_y=B(0,\widehat B)\setminus\overline{B(y,\widehat d)},\qquad
 A=B(-(B+\eta+2)e_n,1/4).
\]
The excluded closed ball lies strictly inside $B(0,\widehat B)$,
so $D_y$ is connected; it contains the fixed lower ball $A$, which
is disjoint from every excluded ball and lies below $x_n=-\eta$.
For fixed $y\in W_\lambda$ in the target range, the constructed kernel
and $G_{1,\lambda}$ solve the same first-variable system on $D_y$.
On $A$ the constructed kernel equals $G_{2,\lambda}$, which equals
$G_{1,\lambda}$ because both variables are in $W_\lambda$.
Unique continuation on $D_y$ proves agreement on the target overlap.
Interchanging the variables proves agreement with the other piece.
The inequalities defining $\widehat B$ and $\widehat d$ leave open
room around the original target, so this also gives the claimed
neighborhood and its compact-subset estimates.
\end{proof}

The constants in Proposition~\ref{cross:initial} may deteriorate as
$d\downarrow0$.  Its role is to initialize the sphere construction at fixed positive
separations.  Lemma~\ref{cross:propagation} will subsequently give
existence at each positive radius of a shrinking family.  The separate
estimate on that family's transfer density, rather than an estimate
obtained by iterating this continuation, will control the shrinking
limit.
